\section{Conclusion}
\label{sec:conclusion}

This paper estimates the fiscal costs of judicial enforcement in public health procurement, exploiting granular bid-level data on pharmaceutical purchases in S\~ao Paulo, Brazil. Court-mandated urgent procurement significantly distorts outcomes: reference prices rise by 2.7\%, negotiated prices increase by 3.1--5.4\%, and bidder participation falls by 4.1--5.5\%, all within item, year, and buyer comparisons.

The ``under the gun'' analysis reveals that the sanction channel is particularly consequential. Comparing litigated and administrative urgent purchases---which share identical planning constraints but differ in penalty exposure---judicial sanctions alone account for a 23--30\% price premium. This magnitude suggests that the threat of judicial punishment fundamentally alters the bargaining environment, distorting procurement incentives well beyond what planning disruption alone would imply. That urgent tenders are more likely to succeed despite worse terms confirms the mechanism: officials under judicial pressure prioritize compliance over cost-efficiency, accepting higher prices to avoid the severe consequences of tender failure. This pattern echoes the broader insight from the accountability literature that oversight mechanisms can generate unintended costs when they distort rather than align incentives \citep{rasul2018management, duflo2018obligations}.

Our results point to a specific, implementable reform. The administrative request mechanism---which since 2009 has allowed S\~ao Paulo to procure medicines for individuals without the threat of judicial sanctions---reduces prices by 23--30\% relative to litigated purchases while delivering the same items under the same planning constraints. Expanding this mechanism statewide and to other Brazilian states could yield substantial fiscal savings. More broadly, institutional designs that decouple the obligation to provide health care from the threat of personal sanctions on procurement officials would preserve accountability while eliminating the ``under the gun'' premium. Complementary supply-side interventions---such as framework agreements or strategic stockpiling for commonly litigated items---could further mitigate enforcement costs. The heterogeneity evidence in the Appendix, showing that the urgency premium is concentrated among basic SUS items and in markets with lower competition, suggests that targeted interventions in these segments could be particularly effective.

The estimated price premiums, applied to the approximately US\$300~million spent annually on litigated health items in S\~ao Paulo alone, imply tens of millions of dollars in excess costs. These findings highlight a tension at the heart of right-to-health enforcement: judicial mandates intended to guarantee individual access impose substantial costs on the public budget, potentially reducing resources available for the broader health system. Institutional designs that balance individual rights with procurement efficiency---rather than treating them as competing objectives---could improve welfare for both litigants and the broader population.
